\documentclass{article}

\usepackage{float}
\usepackage{amsmath}

\begin{document}

\title{COS 314 Assignment Two Report}
% Simple multi-author: each author on its own line
\author{u24594522 -- Keegan Lewis \\
u23588579 -- Tshepiso Makelana \\
u21516261 -- Byron Pillay}
\date{April 2026}
\maketitle

\section{Introduction}

\section{Genetic Algorithm Configuration}

The Genetic Algorithm (GA) was configured with a specific set of the Knapsack Problem. The operators and details are listed below:

\begin{itemize}
    \item \textbf{Termination Criterion:} The algorithm executes for a fixed number of generations, defined by the \texttt{termination} parameter passed to the execution function.
    
    \item \textbf{Population Size:} A constant population of $N = 26$ individuals was maintained throughout the evolutionary process.
    
    \item \textbf{Initial Population and Repair:} The initial population was generated by randomly selecting approximately 25\% of the genes for each chromosome. To ensure feasibility, any individual that exceeded the knapsack capacity was subjected to a \textit{Random Repair} mechanism, where randomly selected items were removed until the total weight was within limits.
    
    \item \textbf{Selection Mechanism:} A hybrid selection strategy was employed:
    \begin{itemize}
        \item \textbf{Elitism:} The single best individual from the current population is preserved and copied directly into the next generation.
        \item \textbf{Roulette Wheel Selection:} The remaining 25 slots are filled using fitness-proportionate selection, where the probability of an individual being chosen was based on their fitness.
    \end{itemize}
    
    \item \textbf{Genetic Operators:}
    \begin{itemize}
        \item \textbf{Crossover:} A \textit{Greedy Intersection Crossover} was applied with a probability of $P_c = 0.85$. This operator preserves items common to both parents and fills the remaining capacity using a greedy approach based on the highest value-to-weight ratios.
        \item \textbf{Mutation:} Bit-flip mutation was applied with a probability of $P_m = 1/n$ (where $n$ is the number of items). 
    \end{itemize}
    
    \item \textbf{Population Update:} After the genetic operators, the population was update in a greedy style where only the best of the parents and children would move on to the next generation.
\end{itemize}

\section{Iterated Local Search Configuration}

The Iterated Local Search (ILS) was configured with a specific set of the Knapsack Problem. The operators and details are listed below:

\begin{itemize}
    \item \textbf{Initial Solution:} The algorithm generates a greedy initial solution by attempting to add each item to the knapsack, unless the item causes the total weight to exceed capacity.
    \item \textbf{Perturbation:} The Perturbation Algorithm takes a solution and randomly applies one of three mutations:
    \begin{itemize}
        \item \textbf 1-for-1 Swap: Removes one random item from the knapsack and adds one random item not in the knapsack to the knapsack.
        \item \textbf 1-for-2 Swap: Removes one random item from the knapsack and adds two random items not in the knapsack to the knapsack.
        \item \textbf 2-for-1 Swap: Removes two random items from the knapsack and adds one random item not in the knapsack to the knapsack.
    \end{itemize}
    If a chosen mutation cannot be performed or produces an invalid solution, the algorithm continues trying until a valid perturbed solution is found or 100 attempts have been made.
    \item \textbf{Termination Condition:} The User specifies the number of iterations the ILS will run. However, ten or more iterations is recommended, and 10 iterations was used during the experiment.
    \item \textbf{Acceptance Criterion:} The current solution is accepted if it achieves a greater fitness than the previous solution. If it is worse , it is still accepted with a probability of 60\%.
    \item \textbf{Local Search:} The Local Search (LS) used by the ILS is a greedy hill climbing algorithm that selects the best neighbour from the current neighbourhood. The neighbourhood consists of all valid solutions that can be obtained by adding one item, removing one item, or swapping one included item with one excluded item.
\end{itemize}

\section{Experimental Setup}

Both algorithms were run on all 11 knapsack problem instances using seed value 2, supplied at runtime via the program's seed prompt. The same seed was applied to both the GA and ILS to ensure reproducibility and to provide a controlled basis for comparison, in accordance with the assignment requirements. The GA was run for 160 generations with a population size of 26. The ILS was run for 100 outer iterations with 70 local search iterations per call.

\subsection*{Wilcoxon Signed-Rank Test}

A one-tailed Wilcoxon signed-rank test was conducted at the 5\% significance level ($\alpha = 0.05$) to assess whether one algorithm produces statistically significantly better solutions than the other. The null hypothesis $H_0$ states that the median difference in best fitness between the two algorithms is zero (i.e., they are equivalent).

The signed differences $d_i = \text{GA}_i - \text{ILS}_i$ were computed for each of the 11 instances. Of these, 10 differences were exactly zero (both algorithms achieved identical fitness), leaving only one non-zero pair: instance f8\_l-d\_kp\_23\_10000, where $d_8 = 9765 - 9767 = -2$ (ILS performed better).

After discarding tied pairs, the effective sample size is $n = 1$. The single absolute difference $|d_8| = 2$ receives rank $R_1 = 1$ with a negative sign, giving $W^{+} = 0$ and $W^{-} = 1$.

For a one-tailed test ($H_1$: ILS $>$ GA), the test statistic is $W^{+} = 0$. With $n = 1$, the minimum achievable $p$-value is $0.5$, which exceeds $\alpha = 0.05$. We therefore \textbf{fail to reject} $H_0$.

\textit{Conclusion:} There is no statistically significant difference between the GA and ILS at the 5\% significance level. Both algorithms achieved the known optimum on 10 out of 11 instances; ILS found the optimum on f8 while the GA fell 2 units short, but this single discrepancy is insufficient to establish a statistically significant advantage.

\section{Results}

\begin{table}[H]
\centering
\caption{Comparison of GA and ILS on 11 knapsack problem instances}
\begin{tabular}{|l|c|c|c|c|c|}
\hline
\textbf{Problem Instance} & \textbf{Algorithm} & \textbf{Seed} & \textbf{Best Solution} & \textbf{Known Optimum} & \textbf{Runtime (s)} \\
\hline
f1\_l-d\_kp\_10\_269 & ILS & 2 & 295.0 & 295 & 0.002 \\
                     & GA  & 2 & 295.0 & 295 & 0.020 \\
\hline
f2\_l-d\_kp\_20\_878 & ILS & 2 & 1024.0 & 1024 & 0.003 \\
                     & GA  & 2 & 1024.0 & 1024 & 0.009 \\
\hline
f3\_l-d\_kp\_4\_20 & ILS & 2 & 35.0 & 35 & 0.000 \\
                   & GA  & 2 & 35.0 & 35 & 0.005 \\
\hline
f4\_l-d\_kp\_4\_11 & ILS & 2 & 23.0 & 23 & 0.000 \\
                                     & GA  & 2 & 23.0 & 23 & 0.005 \\
\hline
f5\_l-d\_kp\_15\_375 & ILS & 2 & 481.0694 & 481.0694 & 0.002 \\
                                         & GA  & 2 & 481.0694 & 481.0694 & 0.009 \\
\hline
f6\_l-d\_kp\_10\_60 & ILS & 2 & 52.0 & 52 & 0.001 \\
                                        & GA  & 2 & 52.0 & 52 & 0.008 \\
\hline
f7\_l-d\_kp\_7\_50 & ILS & 2 & 107.0 & 107 & 0.001 \\
                                     & GA  & 2 & 107.0 & 107 & 0.007 \\
\hline
f8\_l-d\_kp\_23\_10000 & ILS & 2 & 9767.0 & 9767 & 0.005 \\
                       & GA  & 2 & 9765.0 & 9767 & 0.012 \\
\hline
f9\_l-d\_kp\_5\_80 & ILS & 2 & 130.0 & 130 & 0.000 \\
                                     & GA  & 2 & 130.0 & 130 & 0.006 \\
\hline
f10\_l-d\_kp\_20\_879 & ILS & 2 & 1025.0 & 1025 & 0.003 \\
                                            & GA  & 2 & 1025.0 & 1025 & 0.006 \\
\hline
knapPI\_1\_100\_1000\_1 & ILS & 2 & 9147.0 & 9147 & 0.046 \\
                                                & GA  & 2 & 9147.0 & 9147 & 0.021 \\
\hline
\end{tabular}
\end{table}

\end{document}